% Glossary
\selectlanguage{greek}

\newcommand{\gloss}[2]{\item[] \textbf{#1}: #2}

\chapter*{Γλωσσάριο}
\markboth{Γλωσσάριο}{Γλωσσάριο}

\begin{description}
    \gloss{Αλυσίδα Κλιμάκωσης (\en{Cascade})}{Βαθμιδωτή αλυσίδα ταξινομητών όπου κάθε επίπεδο χειρίζεται μόνο τα ζεύγη που δεν έλυσαν τα προηγούμενα, δρομολογώντας τα υπόλοιπα σε ακριβότερο μοντέλο.}
    \gloss{Γράφημα με Κειμενικά Χαρακτηριστικά (\en{Text-Attributed Graph})}{Γράφημα του οποίου οι κόμβοι φέρουν συνοδευτικό κειμενικό περιεχόμενο (π.χ. άρθρα \en{Wikipedia}), επιτρέποντας συνδυασμό δομικού και σημασιολογικού σήματος.}
    \gloss{Δείκτης \en{Adamic-Adar}}{Δομικός ευρετικός δείκτης πρόβλεψης ακμών που σταθμίζει κάθε κοινό γείτονα δύο κόμβων αντιστρόφως ανάλογα με τον λογάριθμο του βαθμού του, δίνοντας μεγαλύτερο βάρος σε σπάνιους κοινούς γείτονες.}
    \gloss{Δείκτες \en{Katz} και \en{SimRank}}{Δομικοί ευρετικοί δείκτες που, σε αντίθεση με τους \en{CN}/\en{Jaccard}/\en{AA}/\en{PA}, αθροίζουν συνεισφορά από διαδρομές αυθαίρετου μήκους μεταξύ δύο κόμβων αντί μόνο από άμεσους κοινούς γείτονες.}
    \gloss{Δρομολόγηση Βάσει Αξιοπιστίας (\en{Confidence routing})}{Μηχανισμός προώθησης ζευγών με αβέβαιη πρόβλεψη (χαμηλή μέγιστη πιθανότητα) στο επόμενο, ακριβότερο επίπεδο της αλυσίδας.}
    \gloss{Δυσκολία Ζεύγους (\en{Difficulty label})}{Ετικέτα ανά ζεύγος κόμβων (\en{trivial\_self\_loop}, \en{trivial\_high\_cn}, \en{trivial\_high\_textsim}, \en{hard}) που χαρακτηρίζει πόσο εύκολα διαχωρίζεται από απλά κριτήρια δομής ή κειμενικής ομοιότητας.}
    \gloss{\en{Cold-start}}{Ζεύγος στο οποίο τουλάχιστον ένα άκρο απουσιάζει από το γράφημα εκπαίδευσης. Στο \en{CascadeLP} αναγνωρίζεται από μηδενικές τιμές και στα τέσσερα δομικά χαρακτηριστικά και παρακάμπτει το Επίπεδο~1.}
    \gloss{Υποσύνολο Μηδενικού \en{CN}}{Διαγνωστικό υποσύνολο ζευγών χωρίς κοινούς γείτονες. Είναι ευρύτερο από το \en{cold-start}, επειδή μπορεί να περιλαμβάνει δύο γνωστούς κόμβους με μη μηδενικούς βαθμούς.}
    \gloss{Ενσωμάτωση (\en{Embedding})}{Χαμηλοδιάστατη διανυσματική αναπαράσταση κόμβου ή κειμένου που καταγράφει σημασιολογική ή δομική ομοιότητα.}
    \gloss{Επιλεκτική Πρόβλεψη (\en{Selective prediction})}{Στρατηγική στην οποία ένα μοντέλο απέχει από την πρόβλεψη (ή την αναθέτει σε ακριβότερο μοντέλο) όταν η αβεβαιότητά του υπερβαίνει ένα όριο.}
    \gloss{\en{Infobox}}{Δομημένο μπλοκ μεταδεδομένων (πεδίο-τιμή) στη σήμανση άρθρου \en{Wikipedia}. Η παρούσα υλοποίηση δεν το αναλύει σε ξεχωριστό πεδίο.}
    \gloss{Καμπύλη \en{ROC}}{Χαρακτηριστική Καμπύλη Λειτουργίας Δέκτη· απεικονίζει τον ρυθμό αληθώς θετικών έναντι του ρυθμού ψευδώς θετικών σε όλα τα όρια απόφασης ενός ταξινομητή.}
    \gloss{Κοινοί Γείτονες (\en{Common Neighbours})}{Δομικός ευρετικός δείκτης πρόβλεψης ακμών που μετρά τον αριθμό των κόμβων που είναι ταυτόχρονα γείτονες και των δύο εξεταζόμενων κόμβων.}
    \gloss{Λανθάνων Χώρος (\en{Latent space})}{Ο χώρος των ενσωματώσεων στον οποίο απεικονίζονται κόμβοι ή κείμενα ώστε η γεωμετρική εγγύτητα να αντιστοιχεί σε σημασιολογική ή δομική ομοιότητα.}
    \gloss{Μεταβίβαση Μηνυμάτων (\en{Message passing})}{Μηχανισμός στα Νευρωνικά Δίκτυα Γραφημάτων όπου η αναπαράσταση κάθε κόμβου ενημερώνεται επαναληπτικά συγκεντρώνοντας πληροφορία από τις αναπαραστάσεις των γειτόνων του.}
    \gloss{Μέρη του Λόγου (\en{Part-of-speech})}{Γραμματική κατηγορία λέξης (ουσιαστικό, ρήμα, επίθετο κ.λπ.)· η κατανομή συχνότητάς τους σε ένα κείμενο χρησιμοποιείται ως χαρακτηριστικό ταξινόμησης.}
    \gloss{Μη Επιβλεπόμενη Μάθηση Αναπαραστάσεων (\en{Representation learning})}{Εκμάθηση χαμηλοδιάστατων αναπαραστάσεων (ενσωματώσεων) κόμβων ή κειμένου χωρίς επιβλεπόμενες ετικέτες, ώστε να αξιοποιηθούν σε επόμενη εργασία.}
    \gloss{\en{Masked Language Modeling} (MLM)}{Στόχος προ-εκπαίδευσης όπου τυχαία \en{tokens} εισόδου αντικαθίστανται με ειδικό σύμβολο κάλυψης και το μοντέλο εκπαιδεύεται να τα προβλέψει από τα συμφραζόμενα.}
    \gloss{Υπερπροσαρμογή (\en{Overfitting})}{Φαινόμενο κατά το οποίο ένα μοντέλο προσαρμόζεται υπερβολικά στα δεδομένα εκπαίδευσης και χάνει ικανότητα γενίκευσης σε αόρατα δεδομένα.}
    \gloss{Ομοφιλία (\en{Homophily})}{Η τάση παρόμοιων κόμβων (π.χ. άρθρων συναφούς θεματολογίας) να συνδέονται μεταξύ τους συχνότερα από τυχαία ζεύγη.}
    \gloss{Κατώφλι Απόφασης (\en{Decision threshold})}{Το όριο πιθανότητας πάνω από το οποίο μια πρόβλεψη ταξινομείται ως θετική· στο \en{CascadeLP} καθορίζει πότε ένα επίπεδο θεωρείται αρκετά σίγουρο ώστε να μην προωθήσει το ζεύγος παρακάτω.}
    \gloss{Πρόβλεψη Ακμών (\en{Link prediction})}{Εργασία συμπερασμού για το εάν υπάρχει (ή θα εμφανιστεί) ακμή μεταξύ δύο κόμβων ενός γραφήματος.}
    \gloss{Πρόβλεψη Σχέσης Εγγράφων (\en{Document relation prediction})}{Στόχος προ-εκπαίδευσης που αντικαθιστά την πρόβλεψη επόμενης πρότασης, ζητώντας από το μοντέλο να προβλέψει τη σχέση (π.χ. συνδεδεμένα μέσω υπερσυνδέσμου) μεταξύ δύο τμημάτων κειμένου.}
    \gloss{Προ-εκπαίδευση (\en{Pre-training})}{Στάδιο εκπαίδευσης ενός μοντέλου σε μεγάλο, γενικό σώμα δεδομένων πριν από την εξειδίκευσή του (\en{fine-tuning}) σε συγκεκριμένη εργασία.}
    \gloss{Προτιμησιακή Σύνδεση (\en{Preferential attachment})}{Δομικός ευρετικός δείκτης που προσεγγίζει την πιθανότητα ακμής ως το γινόμενο των βαθμών των δύο κόμβων, χωρίς να απαιτεί υπολογισμό τομής γειτονιών.}
    \gloss{\en{Siamese} Αρχιτεκτονική (\en{Siamese architecture})}{Νευρωνική αρχιτεκτονική στην οποία το ίδιο μοντέλο κωδικοποιεί ανεξάρτητα δύο εισόδους (π.χ. δύο προτάσεις) σε ενσωματώσεις σταθερού μεγέθους που συγκρίνονται εκ των υστέρων.}
    \gloss{Στρωματοποιημένη Δειγματοληψία (\en{Stratified sampling})}{Δειγματοληψία που διατηρεί την αναλογία των κλάσεων (θετικά/αρνητικά ζεύγη) του πλήρους συνόλου δεδομένων στο υποσύνολο.}
    \gloss{Συντελεστής \en{Jaccard}}{Δομικός ευρετικός δείκτης που κανονικοποιεί τον αριθμό κοινών γειτόνων δύο κόμβων διαιρώντας με το μέγεθος της ένωσης των γειτονιών τους.}
    \gloss{Συχνότητα Όρου -- Αντίστροφη Συχνότητα Εγγράφου (\en{TF-IDF})}{Στατιστική στάθμιση όρων κειμένου που αυξάνει το βάρος όρων συχνών σε ένα έγγραφο αλλά σπάνιων στο σύνολο των εγγράφων· χρησιμοποιείται για ομοιότητα κειμένου μέσω συνημιτόνου.}
    \gloss{Ταξινομητής Τυχαίου Δάσους (\en{Random Forest classifier})}{Σύνολο αποφασιστικών δέντρων εκπαιδευμένων σε τυχαία υποσύνολα δεδομένων και χαρακτηριστικών, των οποίων οι προβλέψεις συνδυάζονται με ψηφοφορία ή μέσο όρο.}
    \gloss{Τριαδικό Κλείσιμο (\en{Triadic closure})}{Η τάση δύο κόμβων με κοινό γείτονα να συνδεθούν τελικά μεταξύ τους, βασική αρχή πίσω από τους δομικούς ευρετικούς δείκτες πρόβλεψης ακμών.}
    \gloss{Τυχαίος Περίπατος (\en{Random walk})}{Ακολουθία κόμβων παραγόμενη με διαδοχική τυχαία μετάβαση σε γειτονικό κόμβο· χρησιμοποιείται για δειγματοληψία γειτονιών ή εκμάθηση αναπαραστάσεων κόμβων.}
    \gloss{Χρόνος Συμπερασμού (\en{Inference latency})}{Ο χρόνος που απαιτείται από ένα εκπαιδευμένο μοντέλο για να παράξει πρόβλεψη. Στις μετρήσεις του \en{CascadeLP} αφορά προϋπολογισμένα χαρακτηριστικά και δεν περιλαμβάνει την εξαγωγή τους.}
\end{description}
